\begin{lemma}[Centre actions on Ext for left modules]
\label{lemma-ext-center-actions-left-modules}
\begin{reference}
\cite[Chapter IV, \S 4.2, p. 188]{Tohoku}
\end{reference}
\begin{history}
The printed source makes the global Ext groups modules over
$Z(\Gamma(X,\mathcal A))$. For a sheaf of noncommutative rings, a globally
central section need not remain central after restriction, so that action is
not defined in general. The canonical scalar ring is instead
$\Gamma(X,\mathcal Z_{\mathcal A})$, where
$\mathcal Z_{\mathcal A}$ is the sheaf-theoretic centre from Modules, Lemma
\ref{modules-lemma-sheaf-hom-left-modules}. The two rings agree when
$\mathcal A$ is commutative.
\end{history}
Let $X$ be a topological space, let $\mathcal A$ be a sheaf of not
necessarily commutative rings with $1$, and let $\mathcal F,\mathcal G$ be
sheaves of unital left $\mathcal A$-modules. Write
$\mathcal Z_{\mathcal A}$ for the sheaf-theoretic centre of $\mathcal A$.
For every $p\geq0$, the global Ext group
\[
\Ext^p_{\mathcal A}(X;\mathcal F,\mathcal G)
=
\Ext^p_{\textit{Mod}(\mathcal A)}(\mathcal F,\mathcal G)
\]
is canonically a $\Gamma(X,\mathcal Z_{\mathcal A})$-module, and the
right-derived sheaf
\[
\SheafExt^p_{\mathcal A}(\mathcal F,\mathcal G)
=
R^p\SheafHom_{\mathcal A}(\mathcal F,-)(\mathcal G)
\]
is canonically a sheaf of $\mathcal Z_{\mathcal A}$-modules. These actions
are functorial in both module arguments and agree in degree zero with the
actions of Modules, Lemma
\ref{modules-lemma-sheaf-hom-left-modules}.

If the canonical inclusion
\[
\Gamma(X,\mathcal Z_{\mathcal A})
\longrightarrow Z(\Gamma(X,\mathcal A))
\]
is an equality---in particular, if $\mathcal A$ is commutative---then the
global Ext groups are modules over $Z(\Gamma(X,\mathcal A))$.
\end{lemma}

\begin{proof}
The category of sheaves of unital left $\mathcal A$-modules is abelian and
has enough injectives; the proof of Injectives, Lemma
\ref{injectives-lemma-sheaves-modules-space} does not use commutativity.
Choose an injective resolution
\[
\mathcal G[0]\longrightarrow\mathcal I^\bullet.
\]
For every $q$, Modules, Lemma
\ref{modules-lemma-sheaf-hom-left-modules} makes
$\SheafHom_{\mathcal A}(\mathcal F,\mathcal I^q)$ a sheaf of
$\mathcal Z_{\mathcal A}$-modules. Since the differentials of
$\mathcal I^\bullet$ are $\mathcal A$-linear, the induced differentials on
these Hom sheaves are $\mathcal Z_{\mathcal A}$-linear. Their cohomology
sheaves are therefore $\mathcal Z_{\mathcal A}$-modules, and they compute
the displayed right-derived sheaves.

Taking global sections gives a complex
\[
\Hom_{\mathcal A}(\mathcal F,\mathcal I^\bullet)
=
\Gamma\bigl(X,
\SheafHom_{\mathcal A}(\mathcal F,\mathcal I^\bullet)\bigr)
\]
of $\Gamma(X,\mathcal Z_{\mathcal A})$-modules. Its degree-$p$ cohomology is
$\Ext^p_{\textit{Mod}(\mathcal A)}(\mathcal F,\mathcal G)$, which gives the
global action. The construction is natural in $\mathcal G$ by functoriality
of derived functors, and natural in $\mathcal F$ by precomposition. Degree
zero is exactly the action constructed in the cited Modules lemma. The final
assertion follows by transporting scalars through the displayed equality.
\end{proof}

